%=============================================================================
%  fig_inductive.tex
%
%  FIGURE 2: Conceptual contrast between transductive (codex) attention and
%  geometric (inductive) attention.
%
%  Layout: two panels, left = codex, right = geometric.
%
%   LEFT  (codex / transductive)         RIGHT  (geometric / inductive)
%   --------------------------           -----------------------------
%   x_i  +  c_i  ->  token               x_i               -> token
%        ^                                                       |
%   electrode-indexed embedding          attention(token_i, token_j, g_ij)
%   added BEFORE attention               geometry enters the attention FUNCTION
%
%   attention is geometry-blind          no per-electrode parameters
%
%  USAGE:    %=============================================================================
%  fig_inductive.tex
%
%  FIGURE 2: Conceptual contrast between transductive (codex) attention and
%  geometric (inductive) attention.
%
%  Layout: two panels, left = codex, right = geometric.
%
%   LEFT  (codex / transductive)         RIGHT  (geometric / inductive)
%   --------------------------           -----------------------------
%   x_i  +  c_i  ->  token               x_i               -> token
%        ^                                                       |
%   electrode-indexed embedding          attention(token_i, token_j, g_ij)
%   added BEFORE attention               geometry enters the attention FUNCTION
%
%   attention is geometry-blind          no per-electrode parameters
%
%  USAGE:    %=============================================================================
%  fig_inductive.tex
%
%  FIGURE 2: Conceptual contrast between transductive (codex) attention and
%  geometric (inductive) attention.
%
%  Layout: two panels, left = codex, right = geometric.
%
%   LEFT  (codex / transductive)         RIGHT  (geometric / inductive)
%   --------------------------           -----------------------------
%   x_i  +  c_i  ->  token               x_i               -> token
%        ^                                                       |
%   electrode-indexed embedding          attention(token_i, token_j, g_ij)
%   added BEFORE attention               geometry enters the attention FUNCTION
%
%   attention is geometry-blind          no per-electrode parameters
%
%  USAGE:    %=============================================================================
%  fig_inductive.tex
%
%  FIGURE 2: Conceptual contrast between transductive (codex) attention and
%  geometric (inductive) attention.
%
%  Layout: two panels, left = codex, right = geometric.
%
%   LEFT  (codex / transductive)         RIGHT  (geometric / inductive)
%   --------------------------           -----------------------------
%   x_i  +  c_i  ->  token               x_i               -> token
%        ^                                                       |
%   electrode-indexed embedding          attention(token_i, token_j, g_ij)
%   added BEFORE attention               geometry enters the attention FUNCTION
%
%   attention is geometry-blind          no per-electrode parameters
%
%  USAGE:    \input{fig_inductive} inside a figure environment.
%            Assumes the same TikZ style header as review_fig_flow.tex.
%            Dashed boxes mark "electrode-specific" entry points.
%
%  DO NOT COMPILE. TikZ only.
%=============================================================================

\begin{figure}[h]
  \centering
  \begin{tikzpicture}[x=1cm, y=1cm]

    % ── Headers ───────────────────────────────────────────────────────────
    \node[amberbox, minimum width=6.6cm, font=\small\bfseries] (hdrL) at (0.0, 0.5)
      {(a) Transductive --- codex-based};
    \node[greenbox, minimum width=6.6cm, font=\small\bfseries] (hdrR) at (8.5, 0.5)
      {(b) Geometric --- proposed};

    % ── Dashed divider ───────────────────────────────────────────────────
    \draw[dashed, gray!60] (4.25, 1.2) -- (4.25, -10.5);

    % ════════════════════════════════════════════════════════════════════
    %  LEFT COLUMN: TRANSDUCTIVE
    % ════════════════════════════════════════════════════════════════════

    % Row 1: patch tokens with electrode LABELS as subtitle
    \node[amberbox, minimum width=2.0cm] (xiL) at (-1.4, -0.8)
      {\small Patch $x_i$\\[-2pt]\scriptsize electrode ``C3''};
    \node[amberbox, minimum width=2.0cm] (xjL) at ( 1.4, -0.8)
      {\small Patch $x_j$\\[-2pt]\scriptsize electrode ``C4''};

    % Row 2: codex lookups keyed by label
    \node[purplebox, minimum width=2.0cm] (ciL) at (-1.4, -2.4)
      {\small $c_i$\\[-2pt]\scriptsize lookup ``C3''};
    \node[purplebox, minimum width=2.0cm] (cjL) at ( 1.4, -2.4)
      {\small $c_j$\\[-2pt]\scriptsize lookup ``C4''};

    \draw[fwd] (xiL) -- (ciL);
    \draw[fwd] (xjL) -- (cjL);

    % Row 3: modified tokens (tilde x = x + c)
    \node[graybox, minimum width=2.0cm] (txiL) at (-1.4, -4.0)
      {\small $\tilde{x}_i$\\[-2pt]\scriptsize $x_i + c_i$};
    \node[graybox, minimum width=2.0cm] (txjL) at ( 1.4, -4.0)
      {\small $\tilde{x}_j$\\[-2pt]\scriptsize $x_j + c_j$};

    \draw[fwd] (ciL) -- (txiL);
    \draw[fwd] (cjL) -- (txjL);

    % Row 4: dot-product attention block
    \node[bluebox, minimum width=4.4cm, minimum height=1.0cm] (attL) at (0.0, -5.7)
      {\small Dot-product attention\\[-2pt]\scriptsize $(\tilde{x}_i)^\top (\tilde{x}_j)/\sqrt{d}$};

    \draw[fwd] (txiL.south) -- (attL.north);
    \draw[fwd] (txjL.south) -- (attL.north);

    % Row 5: output
    \node[amberbox, minimum width=2.6cm] (yL) at (0.0, -7.4)
      {\small Output $y_i$};
    \draw[fwd] (attL) -- (yL);

    % Bottom banner (red): label required
    \node[amberbox, minimum width=6.6cm, font=\scriptsize, align=center,
          minimum height=0.9cm, draw=red!60!black] (banL) at (0.0, -9.0)
      {Requires electrode label\\
       --- not portable across layouts};

    % ════════════════════════════════════════════════════════════════════
    %  RIGHT COLUMN: GEOMETRIC
    % ════════════════════════════════════════════════════════════════════

    % Row 1: patch tokens with COORDINATES as subtitle (no labels)
    \node[greenbox, minimum width=2.0cm] (xiR) at (7.1, -0.8)
      {\small Patch $x_i$\\[-2pt]\scriptsize pos. $p_i \in \mathbb{R}^3$};
    \node[greenbox, minimum width=2.0cm] (xjR) at (9.9, -0.8)
      {\small Patch $x_j$\\[-2pt]\scriptsize pos. $p_j \in \mathbb{R}^3$};

    % Row 2: g_ij computed directly from coordinates (no label step)
    \node[purplebox, minimum width=4.0cm] (gR) at (8.5, -3)
      {\small $g_{ij} \in \mathbb{R}^4$\\[-2pt]
       \scriptsize $[\|p_i - p_j\|,\; p_i - p_j]$};

    \draw[fwd] (xiR.south) -- (gR.north);
    \draw[fwd] (xjR.south) -- (gR.north);

    % Row 3: geometric attention block
    \node[bluebox, minimum width=5.4cm, minimum height=1.0cm] (attR) at (8.5, -5.7)
      {\small Geometric attention (G1)\\[-2pt]
       \scriptsize $a^\top \mathrm{LeakyReLU}(W_a [W_q x_i \,\|\, W_k x_j \,\|\, W_g g_{ij}])$};

    % Three inputs into the attention: x_i, x_j, g_ij
    \draw[fwd] (xiR.west) to[out=180, in=180, looseness=0.5] (attR.west);
    \draw[fwd] (xjR.east) to[out=0, in=0, looseness=0.5]  (attR.east);
    \draw[fwd] (gR)        --                  (attR);

    % Row 5: output
    \node[greenbox, minimum width=2.6cm] (yR) at (8.5, -7.4)
      {\small Output $y_i$};
    \draw[fwd] (attR) -- (yR);

    % Bottom banner (green): no label required
    \node[greenbox, minimum width=6.6cm, font=\scriptsize, align=center,
          minimum height=0.9cm, draw=green!50!black] (banR) at (8.5, -9.0)
      {No electrode label\\
       --- same function on any layout};

  \end{tikzpicture}
  \caption{\textbf{Transductive vs.\ inductive spatial attention.}
  Both panels start from the same EEG patches $x_i, x_j$ and arrive at an
  output $y_i$, but they route electrode information through different
  side channels. (a) The transductive baseline keys a learned codex by
  the electrode \emph{label} (``C3'', ``C4''), retrieving identity
  embeddings $c_i, c_j$ that are added to the patches before a standard
  dot-product attention. The codex table has one row per electrode seen
  during pretraining, so an unseen electrode has no entry. (b) The
  geometric encoder reads electrode \emph{positions} directly from the
  montage, computes the descriptor $g_{ij}$ from them, and feeds it into
  a shared attention function alongside the unmodified patches --- with
  no per-electrode parameters and no label dependency. The same function
  applies to any electrode pair given only its 3-D coordinates,
  including pairs from layouts never seen during pretraining.}
  \label{fig:inductive}
\end{figure}
 inside a figure environment.
%            Assumes the same TikZ style header as review_fig_flow.tex.
%            Dashed boxes mark "electrode-specific" entry points.
%
%  DO NOT COMPILE. TikZ only.
%=============================================================================

\begin{figure}[h]
  \centering
  \begin{tikzpicture}[x=1cm, y=1cm]

    % ── Headers ───────────────────────────────────────────────────────────
    \node[amberbox, minimum width=6.6cm, font=\small\bfseries] (hdrL) at (0.0, 0.5)
      {(a) Transductive --- codex-based};
    \node[greenbox, minimum width=6.6cm, font=\small\bfseries] (hdrR) at (8.5, 0.5)
      {(b) Geometric --- proposed};

    % ── Dashed divider ───────────────────────────────────────────────────
    \draw[dashed, gray!60] (4.25, 1.2) -- (4.25, -10.5);

    % ════════════════════════════════════════════════════════════════════
    %  LEFT COLUMN: TRANSDUCTIVE
    % ════════════════════════════════════════════════════════════════════

    % Row 1: patch tokens with electrode LABELS as subtitle
    \node[amberbox, minimum width=2.0cm] (xiL) at (-1.4, -0.8)
      {\small Patch $x_i$\\[-2pt]\scriptsize electrode ``C3''};
    \node[amberbox, minimum width=2.0cm] (xjL) at ( 1.4, -0.8)
      {\small Patch $x_j$\\[-2pt]\scriptsize electrode ``C4''};

    % Row 2: codex lookups keyed by label
    \node[purplebox, minimum width=2.0cm] (ciL) at (-1.4, -2.4)
      {\small $c_i$\\[-2pt]\scriptsize lookup ``C3''};
    \node[purplebox, minimum width=2.0cm] (cjL) at ( 1.4, -2.4)
      {\small $c_j$\\[-2pt]\scriptsize lookup ``C4''};

    \draw[fwd] (xiL) -- (ciL);
    \draw[fwd] (xjL) -- (cjL);

    % Row 3: modified tokens (tilde x = x + c)
    \node[graybox, minimum width=2.0cm] (txiL) at (-1.4, -4.0)
      {\small $\tilde{x}_i$\\[-2pt]\scriptsize $x_i + c_i$};
    \node[graybox, minimum width=2.0cm] (txjL) at ( 1.4, -4.0)
      {\small $\tilde{x}_j$\\[-2pt]\scriptsize $x_j + c_j$};

    \draw[fwd] (ciL) -- (txiL);
    \draw[fwd] (cjL) -- (txjL);

    % Row 4: dot-product attention block
    \node[bluebox, minimum width=4.4cm, minimum height=1.0cm] (attL) at (0.0, -5.7)
      {\small Dot-product attention\\[-2pt]\scriptsize $(\tilde{x}_i)^\top (\tilde{x}_j)/\sqrt{d}$};

    \draw[fwd] (txiL.south) -- (attL.north);
    \draw[fwd] (txjL.south) -- (attL.north);

    % Row 5: output
    \node[amberbox, minimum width=2.6cm] (yL) at (0.0, -7.4)
      {\small Output $y_i$};
    \draw[fwd] (attL) -- (yL);

    % Bottom banner (red): label required
    \node[amberbox, minimum width=6.6cm, font=\scriptsize, align=center,
          minimum height=0.9cm, draw=red!60!black] (banL) at (0.0, -9.0)
      {Requires electrode label\\
       --- not portable across layouts};

    % ════════════════════════════════════════════════════════════════════
    %  RIGHT COLUMN: GEOMETRIC
    % ════════════════════════════════════════════════════════════════════

    % Row 1: patch tokens with COORDINATES as subtitle (no labels)
    \node[greenbox, minimum width=2.0cm] (xiR) at (7.1, -0.8)
      {\small Patch $x_i$\\[-2pt]\scriptsize pos. $p_i \in \mathbb{R}^3$};
    \node[greenbox, minimum width=2.0cm] (xjR) at (9.9, -0.8)
      {\small Patch $x_j$\\[-2pt]\scriptsize pos. $p_j \in \mathbb{R}^3$};

    % Row 2: g_ij computed directly from coordinates (no label step)
    \node[purplebox, minimum width=4.0cm] (gR) at (8.5, -3)
      {\small $g_{ij} \in \mathbb{R}^4$\\[-2pt]
       \scriptsize $[\|p_i - p_j\|,\; p_i - p_j]$};

    \draw[fwd] (xiR.south) -- (gR.north);
    \draw[fwd] (xjR.south) -- (gR.north);

    % Row 3: geometric attention block
    \node[bluebox, minimum width=5.4cm, minimum height=1.0cm] (attR) at (8.5, -5.7)
      {\small Geometric attention (G1)\\[-2pt]
       \scriptsize $a^\top \mathrm{LeakyReLU}(W_a [W_q x_i \,\|\, W_k x_j \,\|\, W_g g_{ij}])$};

    % Three inputs into the attention: x_i, x_j, g_ij
    \draw[fwd] (xiR.west) to[out=180, in=180, looseness=0.5] (attR.west);
    \draw[fwd] (xjR.east) to[out=0, in=0, looseness=0.5]  (attR.east);
    \draw[fwd] (gR)        --                  (attR);

    % Row 5: output
    \node[greenbox, minimum width=2.6cm] (yR) at (8.5, -7.4)
      {\small Output $y_i$};
    \draw[fwd] (attR) -- (yR);

    % Bottom banner (green): no label required
    \node[greenbox, minimum width=6.6cm, font=\scriptsize, align=center,
          minimum height=0.9cm, draw=green!50!black] (banR) at (8.5, -9.0)
      {No electrode label\\
       --- same function on any layout};

  \end{tikzpicture}
  \caption{\textbf{Transductive vs.\ inductive spatial attention.}
  Both panels start from the same EEG patches $x_i, x_j$ and arrive at an
  output $y_i$, but they route electrode information through different
  side channels. (a) The transductive baseline keys a learned codex by
  the electrode \emph{label} (``C3'', ``C4''), retrieving identity
  embeddings $c_i, c_j$ that are added to the patches before a standard
  dot-product attention. The codex table has one row per electrode seen
  during pretraining, so an unseen electrode has no entry. (b) The
  geometric encoder reads electrode \emph{positions} directly from the
  montage, computes the descriptor $g_{ij}$ from them, and feeds it into
  a shared attention function alongside the unmodified patches --- with
  no per-electrode parameters and no label dependency. The same function
  applies to any electrode pair given only its 3-D coordinates,
  including pairs from layouts never seen during pretraining.}
  \label{fig:inductive}
\end{figure}
 inside a figure environment.
%            Assumes the same TikZ style header as review_fig_flow.tex.
%            Dashed boxes mark "electrode-specific" entry points.
%
%  DO NOT COMPILE. TikZ only.
%=============================================================================

\begin{figure}[h]
  \centering
  \begin{tikzpicture}[x=1cm, y=1cm]

    % ── Headers ───────────────────────────────────────────────────────────
    \node[amberbox, minimum width=6.6cm, font=\small\bfseries] (hdrL) at (0.0, 0.5)
      {(a) Transductive --- codex-based};
    \node[greenbox, minimum width=6.6cm, font=\small\bfseries] (hdrR) at (8.5, 0.5)
      {(b) Geometric --- proposed};

    % ── Dashed divider ───────────────────────────────────────────────────
    \draw[dashed, gray!60] (4.25, 1.2) -- (4.25, -10.5);

    % ════════════════════════════════════════════════════════════════════
    %  LEFT COLUMN: TRANSDUCTIVE
    % ════════════════════════════════════════════════════════════════════

    % Row 1: patch tokens with electrode LABELS as subtitle
    \node[amberbox, minimum width=2.0cm] (xiL) at (-1.4, -0.8)
      {\small Patch $x_i$\\[-2pt]\scriptsize electrode ``C3''};
    \node[amberbox, minimum width=2.0cm] (xjL) at ( 1.4, -0.8)
      {\small Patch $x_j$\\[-2pt]\scriptsize electrode ``C4''};

    % Row 2: codex lookups keyed by label
    \node[purplebox, minimum width=2.0cm] (ciL) at (-1.4, -2.4)
      {\small $c_i$\\[-2pt]\scriptsize lookup ``C3''};
    \node[purplebox, minimum width=2.0cm] (cjL) at ( 1.4, -2.4)
      {\small $c_j$\\[-2pt]\scriptsize lookup ``C4''};

    \draw[fwd] (xiL) -- (ciL);
    \draw[fwd] (xjL) -- (cjL);

    % Row 3: modified tokens (tilde x = x + c)
    \node[graybox, minimum width=2.0cm] (txiL) at (-1.4, -4.0)
      {\small $\tilde{x}_i$\\[-2pt]\scriptsize $x_i + c_i$};
    \node[graybox, minimum width=2.0cm] (txjL) at ( 1.4, -4.0)
      {\small $\tilde{x}_j$\\[-2pt]\scriptsize $x_j + c_j$};

    \draw[fwd] (ciL) -- (txiL);
    \draw[fwd] (cjL) -- (txjL);

    % Row 4: dot-product attention block
    \node[bluebox, minimum width=4.4cm, minimum height=1.0cm] (attL) at (0.0, -5.7)
      {\small Dot-product attention\\[-2pt]\scriptsize $(\tilde{x}_i)^\top (\tilde{x}_j)/\sqrt{d}$};

    \draw[fwd] (txiL.south) -- (attL.north);
    \draw[fwd] (txjL.south) -- (attL.north);

    % Row 5: output
    \node[amberbox, minimum width=2.6cm] (yL) at (0.0, -7.4)
      {\small Output $y_i$};
    \draw[fwd] (attL) -- (yL);

    % Bottom banner (red): label required
    \node[amberbox, minimum width=6.6cm, font=\scriptsize, align=center,
          minimum height=0.9cm, draw=red!60!black] (banL) at (0.0, -9.0)
      {Requires electrode label\\
       --- not portable across layouts};

    % ════════════════════════════════════════════════════════════════════
    %  RIGHT COLUMN: GEOMETRIC
    % ════════════════════════════════════════════════════════════════════

    % Row 1: patch tokens with COORDINATES as subtitle (no labels)
    \node[greenbox, minimum width=2.0cm] (xiR) at (7.1, -0.8)
      {\small Patch $x_i$\\[-2pt]\scriptsize pos. $p_i \in \mathbb{R}^3$};
    \node[greenbox, minimum width=2.0cm] (xjR) at (9.9, -0.8)
      {\small Patch $x_j$\\[-2pt]\scriptsize pos. $p_j \in \mathbb{R}^3$};

    % Row 2: g_ij computed directly from coordinates (no label step)
    \node[purplebox, minimum width=4.0cm] (gR) at (8.5, -3)
      {\small $g_{ij} \in \mathbb{R}^4$\\[-2pt]
       \scriptsize $[\|p_i - p_j\|,\; p_i - p_j]$};

    \draw[fwd] (xiR.south) -- (gR.north);
    \draw[fwd] (xjR.south) -- (gR.north);

    % Row 3: geometric attention block
    \node[bluebox, minimum width=5.4cm, minimum height=1.0cm] (attR) at (8.5, -5.7)
      {\small Geometric attention (G1)\\[-2pt]
       \scriptsize $a^\top \mathrm{LeakyReLU}(W_a [W_q x_i \,\|\, W_k x_j \,\|\, W_g g_{ij}])$};

    % Three inputs into the attention: x_i, x_j, g_ij
    \draw[fwd] (xiR.west) to[out=180, in=180, looseness=0.5] (attR.west);
    \draw[fwd] (xjR.east) to[out=0, in=0, looseness=0.5]  (attR.east);
    \draw[fwd] (gR)        --                  (attR);

    % Row 5: output
    \node[greenbox, minimum width=2.6cm] (yR) at (8.5, -7.4)
      {\small Output $y_i$};
    \draw[fwd] (attR) -- (yR);

    % Bottom banner (green): no label required
    \node[greenbox, minimum width=6.6cm, font=\scriptsize, align=center,
          minimum height=0.9cm, draw=green!50!black] (banR) at (8.5, -9.0)
      {No electrode label\\
       --- same function on any layout};

  \end{tikzpicture}
  \caption{\textbf{Transductive vs.\ inductive spatial attention.}
  Both panels start from the same EEG patches $x_i, x_j$ and arrive at an
  output $y_i$, but they route electrode information through different
  side channels. (a) The transductive baseline keys a learned codex by
  the electrode \emph{label} (``C3'', ``C4''), retrieving identity
  embeddings $c_i, c_j$ that are added to the patches before a standard
  dot-product attention. The codex table has one row per electrode seen
  during pretraining, so an unseen electrode has no entry. (b) The
  geometric encoder reads electrode \emph{positions} directly from the
  montage, computes the descriptor $g_{ij}$ from them, and feeds it into
  a shared attention function alongside the unmodified patches --- with
  no per-electrode parameters and no label dependency. The same function
  applies to any electrode pair given only its 3-D coordinates,
  including pairs from layouts never seen during pretraining.}
  \label{fig:inductive}
\end{figure}
 inside a figure environment.
%            Assumes the same TikZ style header as review_fig_flow.tex.
%            Dashed boxes mark "electrode-specific" entry points.
%
%  DO NOT COMPILE. TikZ only.
%=============================================================================

\begin{figure}[h]
  \centering
  \begin{tikzpicture}[x=1cm, y=1cm]

    % ── Headers ───────────────────────────────────────────────────────────
    \node[amberbox, minimum width=6.6cm, font=\small\bfseries] (hdrL) at (0.0, 0.5)
      {(a) Transductive --- codex-based};
    \node[greenbox, minimum width=6.6cm, font=\small\bfseries] (hdrR) at (8.5, 0.5)
      {(b) Geometric --- proposed};

    % ── Dashed divider ───────────────────────────────────────────────────
    \draw[dashed, gray!60] (4.25, 1.2) -- (4.25, -10.5);

    % ════════════════════════════════════════════════════════════════════
    %  LEFT COLUMN: TRANSDUCTIVE
    % ════════════════════════════════════════════════════════════════════

    % Row 1: patch tokens with electrode LABELS as subtitle
    \node[amberbox, minimum width=2.0cm] (xiL) at (-1.4, -0.8)
      {\small Patch $x_i$\\[-2pt]\scriptsize electrode ``C3''};
    \node[amberbox, minimum width=2.0cm] (xjL) at ( 1.4, -0.8)
      {\small Patch $x_j$\\[-2pt]\scriptsize electrode ``C4''};

    % Row 2: codex lookups keyed by label
    \node[purplebox, minimum width=2.0cm] (ciL) at (-1.4, -2.4)
      {\small $c_i$\\[-2pt]\scriptsize lookup ``C3''};
    \node[purplebox, minimum width=2.0cm] (cjL) at ( 1.4, -2.4)
      {\small $c_j$\\[-2pt]\scriptsize lookup ``C4''};

    \draw[fwd] (xiL) -- (ciL);
    \draw[fwd] (xjL) -- (cjL);

    % Row 3: modified tokens (tilde x = x + c)
    \node[graybox, minimum width=2.0cm] (txiL) at (-1.4, -4.0)
      {\small $\tilde{x}_i$\\[-2pt]\scriptsize $x_i + c_i$};
    \node[graybox, minimum width=2.0cm] (txjL) at ( 1.4, -4.0)
      {\small $\tilde{x}_j$\\[-2pt]\scriptsize $x_j + c_j$};

    \draw[fwd] (ciL) -- (txiL);
    \draw[fwd] (cjL) -- (txjL);

    % Row 4: dot-product attention block
    \node[bluebox, minimum width=4.4cm, minimum height=1.0cm] (attL) at (0.0, -5.7)
      {\small Dot-product attention\\[-2pt]\scriptsize $(\tilde{x}_i)^\top (\tilde{x}_j)/\sqrt{d}$};

    \draw[fwd] (txiL.south) -- (attL.north);
    \draw[fwd] (txjL.south) -- (attL.north);

    % Row 5: output
    \node[amberbox, minimum width=2.6cm] (yL) at (0.0, -7.4)
      {\small Output $y_i$};
    \draw[fwd] (attL) -- (yL);

    % Bottom banner (red): label required
    \node[amberbox, minimum width=6.6cm, font=\scriptsize, align=center,
          minimum height=0.9cm, draw=red!60!black] (banL) at (0.0, -9.0)
      {Requires electrode label\\
       --- not portable across layouts};

    % ════════════════════════════════════════════════════════════════════
    %  RIGHT COLUMN: GEOMETRIC
    % ════════════════════════════════════════════════════════════════════

    % Row 1: patch tokens with COORDINATES as subtitle (no labels)
    \node[greenbox, minimum width=2.0cm] (xiR) at (7.1, -0.8)
      {\small Patch $x_i$\\[-2pt]\scriptsize pos. $p_i \in \mathbb{R}^3$};
    \node[greenbox, minimum width=2.0cm] (xjR) at (9.9, -0.8)
      {\small Patch $x_j$\\[-2pt]\scriptsize pos. $p_j \in \mathbb{R}^3$};

    % Row 2: g_ij computed directly from coordinates (no label step)
    \node[purplebox, minimum width=4.0cm] (gR) at (8.5, -3)
      {\small $g_{ij} \in \mathbb{R}^4$\\[-2pt]
       \scriptsize $[\|p_i - p_j\|,\; p_i - p_j]$};

    \draw[fwd] (xiR.south) -- (gR.north);
    \draw[fwd] (xjR.south) -- (gR.north);

    % Row 3: geometric attention block
    \node[bluebox, minimum width=5.4cm, minimum height=1.0cm] (attR) at (8.5, -5.7)
      {\small Geometric attention (G1)\\[-2pt]
       \scriptsize $a^\top \mathrm{LeakyReLU}(W_a [W_q x_i \,\|\, W_k x_j \,\|\, W_g g_{ij}])$};

    % Three inputs into the attention: x_i, x_j, g_ij
    \draw[fwd] (xiR.west) to[out=180, in=180, looseness=0.5] (attR.west);
    \draw[fwd] (xjR.east) to[out=0, in=0, looseness=0.5]  (attR.east);
    \draw[fwd] (gR)        --                  (attR);

    % Row 5: output
    \node[greenbox, minimum width=2.6cm] (yR) at (8.5, -7.4)
      {\small Output $y_i$};
    \draw[fwd] (attR) -- (yR);

    % Bottom banner (green): no label required
    \node[greenbox, minimum width=6.6cm, font=\scriptsize, align=center,
          minimum height=0.9cm, draw=green!50!black] (banR) at (8.5, -9.0)
      {No electrode label\\
       --- same function on any layout};

  \end{tikzpicture}
  \caption{\textbf{Transductive vs.\ inductive spatial attention.}
  Both panels start from the same EEG patches $x_i, x_j$ and arrive at an
  output $y_i$, but they route electrode information through different
  side channels. (a) The transductive baseline keys a learned codex by
  the electrode \emph{label} (``C3'', ``C4''), retrieving identity
  embeddings $c_i, c_j$ that are added to the patches before a standard
  dot-product attention. The codex table has one row per electrode seen
  during pretraining, so an unseen electrode has no entry. (b) The
  geometric encoder reads electrode \emph{positions} directly from the
  montage, computes the descriptor $g_{ij}$ from them, and feeds it into
  a shared attention function alongside the unmodified patches --- with
  no per-electrode parameters and no label dependency. The same function
  applies to any electrode pair given only its 3-D coordinates,
  including pairs from layouts never seen during pretraining.}
  \label{fig:inductive}
\end{figure}
